\begin{table}[H]
\centering
\caption{Synthetic difference-in-differences estimates}
\label{tab:v1_sdid}
\begin{threeparttable}
\small
\setlength{\tabcolsep}{4pt}
\begin{tabular}{lcccc}
\toprule
& \multicolumn{2}{c}{\# of registered unemployed} & \multicolumn{2}{c}{\# of new contracts} \\
\cmidrule(lr){2-3}\cmidrule(lr){4-5}
& (1) & (2) & (3) & (4) \\
\midrule
High exposure: adjustment period & 0.027$^{**}$ & 0.029$^{**}$ & -0.000 & 0.024 \\
 & (0.011) & (0.013) & (0.022) & (0.030) \\
Impact (percent) & 2.7 & 2.9 & -0.0 & 2.5 \\
\addlinespace
High exposure: later period & 0.041$^{*}$ & 0.019 & 0.008 & 0.029 \\
 & (0.024) & (0.030) & (0.029) & (0.037) \\
Impact (percent) & 4.1 & 1.9 & 0.8 & 3.0 \\
\midrule
Exposure threshold for treatment status & 0.1169 & 0.2 & 0.1169 & 0.2 \\
All lower-exposure occupations eligible as donors & Yes & No & Yes & No \\
Donor pool restricted to zero exposure & No & Yes & No & Yes \\
CNO1 $\times$ month residualization & Yes & Yes & Yes & Yes \\
Treated occupations: adjustment period & 124 & 84 & 118 & 82 \\
Donor occupations: adjustment period & 377 & 252 & 365 & 246 \\
Observations: adjustment period & 23,547 & 15,792 & 22,701 & 15,416 \\
Treated occupations: later period & 124 & 84 & 118 & 82 \\
Donor occupations: later period & 377 & 252 & 363 & 245 \\
Observations: later period & 19,038 & 12,768 & 18,278 & 12,426 \\
\bottomrule
\end{tabular}
\begin{tablenotes}[flushleft]
\footnotesize
\item \emph{Notes:} Columns 1 and 3 report the baseline SDID design, where treated occupations have nearest-neighbor exposure above 0.1169 and every occupation at or below that cutoff remains eligible for the donor pool. Columns 2 and 4 report the more dichotomous robustness design, where treated occupations have exposure above 0.2 and donors are restricted to zero-exposure occupations. The adjustment-period estimates use all pre-treatment months and event times 0--24. The later-period estimates use all pre-treatment months and event times 25--40; event times 0--24 are omitted from those fits. Unit and time weights are re-estimated separately for each specification, outcome, and interval. All columns residualize outcomes on CNO1-by-month indicators before constructing the synthetic comparison. Because this residualization depends only on outcomes and family-month cells, it is held fixed across placebo assignments. Standard errors in parentheses use placebo inference with 500 repetitions. Impact rows report $100[\exp(\widehat{\tau})-1]$. $^{***}p<0.01$, $^{**}p<0.05$, and $^{*}p<0.10$.
\end{tablenotes}
\end{threeparttable}
\end{table}
